\documentclass[11pt,a4paper]{article}
\usepackage[T1]{fontenc}
\usepackage[utf8]{inputenc}
\usepackage{lmodern}
\usepackage{microtype}
\usepackage{amsmath,amssymb,amsthm,mathtools}
\usepackage{geometry}
\usepackage{enumitem}
\usepackage{booktabs,tabularx,array}
\usepackage{xcolor}
\usepackage{hyperref}
\usepackage{fancyhdr}
\usepackage{titlesec}
\usepackage{needspace}
\usepackage{enumitem}
\geometry{margin=29mm,headheight=15pt,headsep=14pt}
\definecolor{ink}{HTML}{1E3045}
\definecolor{accent}{HTML}{235B84}
\definecolor{quiet}{HTML}{5E7182}
\hypersetup{colorlinks=true,linkcolor=accent,urlcolor=accent,citecolor=accent,pdftitle={Local Sobolev Continuity and Finite-Order Parabolic Smoothing},pdfauthor={Mathematical exposition}}
\titleformat{\section}{\normalfont\Large\bfseries\color{ink}}{\thesection}{0.8em}{}
\titleformat{\subsection}{\normalfont\large\bfseries\color{ink}}{\thesubsection}{0.8em}{}
\titleformat{\subsubsection}{\normalfont\normalsize\bfseries\color{ink}}{\thesubsubsection}{0.8em}{}
\titlespacing*{\section}{0pt}{18pt}{7pt}
\titlespacing*{\subsection}{0pt}{13pt}{5pt}
\pagestyle{fancy}
\fancyhf{}
\fancyhead[L]{\small\color{quiet}Local parabolic Sobolev regularity}
\fancyhead[R]{\small\color{quiet}Non-divergence-form operators}
\fancyfoot[C]{\small\thepage}
\renewcommand{\headrulewidth}{0.3pt}
\setlength{\parindent}{0pt}
\setlength{\parskip}{5.2pt}
\setlist[itemize]{topsep=3pt,itemsep=2pt,leftmargin=1.6em}
\setlist[enumerate]{topsep=3pt,itemsep=2pt,leftmargin=1.8em}
\allowdisplaybreaks[1]
\newtheorem{theorem}{Theorem}[section]
\newtheorem{lemma}[theorem]{Lemma}
\newtheorem{corollary}[theorem]{Corollary}
\newtheorem{proposition}[theorem]{Proposition}
\newtheorem{example}[theorem]{Example}
\theoremstyle{remark}
\newtheorem{remark}[theorem]{Remark}
\newtheorem{caution}[theorem]{Technical clarification}
\newcommand{\loc}{\mathrm{loc}}
\newcommand{\inter}{\operatorname{int}}
\newcommand{\dd}{\mathrm{d}}
\newcommand{\R}{\mathbb R}
\newcommand{\N}{\mathbb N}
\newcommand{\cD}{\mathcal D}
\newcommand{\supp}{\operatorname{supp}}
\newcommand{\norm}[1]{\left\lVert#1\right\rVert}
\newcommand{\abs}[1]{\left\lvert#1\right\rvert}
\newcommand{\br}[1]{\left(#1\right)}
\newcommand{\bluebox}[1]{\begin{center}\fcolorbox{accent!55}{accent!4}{\begin{minipage}{0.91\linewidth}\centering\strut #1\strut\end{minipage}}\end{center}}

\begin{document}
\begin{titlepage}
\vspace*{34mm}
{\color{accent}\rule{\linewidth}{1.3pt}}\par\vspace{12mm}
{\Huge\bfseries\color{ink}Local Sobolev Continuity\\[7pt]
and Finite-Order\\[7pt]
Parabolic Smoothing\par}
\vspace{9mm}
{\Large\color{quiet}Initial-time behavior for non-divergence-form\\
elliptic operators with $C^{k,1}$ coefficients\par}
\vspace{15mm}
{\color{accent}\rule{\linewidth}{0.7pt}}\par
\vspace{15mm}
{\large A theorem--proof exposition with energy estimates, elliptic bootstrapping, examples, counterexamples, and derivative corollaries\par}
\vfill
\begin{minipage}{\linewidth}
\small\color{quiet}
\textbf{Background references.} M. E. Taylor, \emph{Partial Differential Equations I: Basic Theory}, 3rd ed., Springer (2023), and N. V. Krylov, \emph{Lectures on Elliptic and Parabolic Equations in Sobolev Spaces}, AMS (2008).\par
\vspace{7pt}
\textbf{Scope.} Local interior regularity and convergence at $t=0$ for solutions satisfying the natural local energy and initial-trace hypotheses. Neither completeness nor a prescribed boundary condition is used in the interior arguments.
\end{minipage}
\vspace{20mm}
\end{titlepage}
\tableofcontents
\newpage

\section{Purpose and the three questions}

Let $M$ be a smooth Riemannian manifold, possibly noncompact and possibly with boundary. Given a time-independent second-order, non-divergence-form elliptic operator $P$ with $C^{k,1}$ coefficients, we wish to analyze the heat-type equation
\begin{equation}\label{eq:heat}
\partial_tu+Pu=0,\qquad u(0)=u_0.
\end{equation}
The objective is \emph{local in space}, including behavior as $t\downarrow0$. Three separate issues arise:
\begin{enumerate}
\item For which $s\geq0$ does $u_0\in H^s_{\loc}$ imply strong convergence $u(t)\to u_0$ in $H^s_{\loc}$?
\item Given the limited coefficient regularity, what is the largest guaranteed spatial Sobolev order at each positive time?
\item For a stronger order $r$, what is the behavior of $\norm{u(t)}_{H^r(U)}$ as $t\downarrow0$?
\end{enumerate}
A local regularity theorem answers all three without constructing a heat semigroup for each possible boundary condition. The main distinction is between the \emph{local regularity of an already existing solution} and the \emph{global well-posedness problem} that produces that solution.

\section{The consolidated theorem}

Let $M^\circ=M\setminus\partial M$. In local coordinates consider
\begin{equation}\label{eq:operator}
P=-a^{ij}(x)\nabla_i\nabla_j+b^i(x)\nabla_i+c(x),
\end{equation}
where $a^{ij},b^i,c\in C^{k,1}_{\loc}(M^\circ)$, $k\in\N_0$. Assume the scalar real-coefficient principal part is strongly elliptic on compact subsets: for each $K\Subset M^\circ$ there exists $\lambda_K>0$ with
\begin{equation}\label{eq:elliptic}
a^{ij}(x)\xi_i\xi_j\geq\lambda_K\abs{\xi}_g^2,
\qquad x\in K.
\end{equation}
The coefficients do not depend on $t$. A \emph{local energy solution up to $t=0$} means a distributional solution of \eqref{eq:heat} such that, for every $V\Subset M^\circ$ and $T'<T$,
\begin{equation}\label{eq:energyclass}
u\in L^2(0,T';H^1(V)),\quad
u_t\in L^2(0,T';H^{-1}(V)),\quad
u(0)=u_0\in L^2_{\loc}(M^\circ)
\end{equation}
with the initial equality interpreted via the $L^2$-continuous representative. An equivalent local formulation uses smooth cutoffs $\chi u$ on coordinate balls; this avoids requiring smooth boundaries for arbitrary $V$.

\begin{theorem}[Local Sobolev continuity and finite-order smoothing]\label{thm:main}
Set
\begin{equation}\label{eq:R}
R=k+3.
\end{equation}
Under \eqref{eq:operator}--\eqref{eq:energyclass}, the following hold.

\smallskip
\textbf{(A) Initial-time strong continuity.} For every real $0\leq s\leq R$,
\begin{equation}\label{eq:A}
u_0\in H^s_{\loc}(M^\circ)
\quad\Longrightarrow\quad
u\in C\bigl([0,T);H^s_{\loc}(M^\circ)\bigr).
\end{equation}
In particular, for each $U\Subset M^\circ$,
\begin{equation}\label{eq:strong}
\lim_{t\downarrow0}\norm{u(t)-u_0}_{H^s(U)}=0.
\end{equation}

\smallskip
\textbf{(B) Positive-time smoothing.} Even when only $u_0\in L^2_{\loc}$ is assumed,
\begin{equation}\label{eq:B}
u\in C\bigl((0,T);H^R_{\loc}(M^\circ)\bigr).
\end{equation}
Thus $u(t)\in H^{k+3}_{\loc}$ for every $t>0$. The integer threshold $k+3$ is sharp as a coefficient-dependent guarantee.

\smallskip
\textbf{(C) Limits of stronger norms.} Fix $U\Subset M^\circ$ and $0\leq r\leq R$.
\begin{itemize}
\item If $u_0\in H^r(V)$ for some $U\Subset V\Subset M^\circ$, then
\begin{equation}\label{eq:finite}
\lim_{t\downarrow0}\norm{u(t)}_{H^r(U)}=\norm{u_0}_{H^r(U)}.
\end{equation}
\item If $u_0\notin H^r(U)$, then
\begin{equation}\label{eq:infinite}
\lim_{t\downarrow0}\norm{u(t)}_{H^r(U)}=+\infty.
\end{equation}
\end{itemize}
When $u_0\in H^r(U)$ but has no $H^r$ regularity on a neighborhood of $\overline U$, neither alternative is automatic.

\smallskip
\textbf{(D) Independence of the boundary realization.} The conclusions apply to every boundary realization for which the corresponding solution satisfies \eqref{eq:energyclass}; the local constants do not involve the boundary operator. Numerical estimates may nevertheless depend on lower-order norms of the particular solution.
\end{theorem}

\begin{remark}[Scope of the proof]
The integer-order estimates below are developed explicitly using energy arguments, commutators, and the elliptic estimate. Extending the estimates to every noninteger exponent requires standard fractional Sobolev commutator/localization estimates or an equivalent parabolic parametrix. We state that analytic input precisely in \S\ref{sec:fractional}, rather than presenting the elementary differentiation calculation as a complete substitute for it. Likewise, passing from formal differentiated equations to weak solutions uses standard difference-quotient regularization and the initial-trace theorem; see \cite{Taylor2023,Krylov2008}.
\end{remark}

\section{The basic local parabolic estimate}\label{sec:energy}

We work in one coordinate chart; the geometry contributes only smooth coefficient terms. Replacing $a^{ij}$ by its symmetric part does not alter the principal second-derivative term. In Euclidean notation,
\begin{equation}\label{eq:localP}
P=-a^{ij}(x)\partial_i\partial_j+b^i(x)\partial_i+c(x).
\end{equation}

\begin{lemma}[Local energy estimate and continuity]\label{lem:energy}
Let $q$ be an integer with $0\leq q\leq k+1$. If $v$ is a local energy solution of $v_t+Pv=0$ and $v_0\in H^q_{\loc}$, then, for $U\Subset V\Subset M^\circ$ and $0<T'<T$,
\begin{equation}\label{eq:energy}
\begin{aligned}
&\sup_{0\leq t\leq T'}\norm{v(t)}_{H^q(U)}^2
+\int_0^{T'}\norm{v(t)}_{H^{q+1}(U)}^2\,\dd t\\
&\qquad\leq C\left(\norm{v_0}_{H^q(V)}^2
+\int_0^T\norm{v(t)}_{L^2(V)}^2\,\dd t\right).
\end{aligned}
\end{equation}
Moreover $v\in C([0,T);H^q_{\loc})$. There is an analogous positive-time estimate, without $H^q$ initial regularity, on every cylinder compactly contained in $M^\circ\times(0,T)$.
\end{lemma}

\subsection{The case $q=0$}
Choose $\chi\in C_c^\infty(V)$ with $\chi=1$ on $U$. Formally multiply the equation
\[
v_t-a^{ij}\partial_i\partial_jv+b^i\partial_iv+cv=0
\]
by $\chi^2v$, and integrate in space. Integration by parts gives
\begin{equation}\label{eq:ibp}
\begin{aligned}
-\int\chi^2v\,a^{ij}\partial_i\partial_jv
={}&\int\chi^2a^{ij}\partial_iv\,\partial_jv\\
&+2\int\chi(\partial_i\chi)a^{ij}v\,\partial_jv\\
&+\int\chi^2(\partial_i a^{ij})v\,\partial_jv.
\end{aligned}
\end{equation}
Ellipticity controls the first term from below by $\lambda\norm{\chi\nabla v}_{L^2}^2$. The remaining terms and the first- and zeroth-order contributions are controlled by Young's inequality. Consequently,
\begin{equation}\label{eq:baseenergy}
\frac{\dd}{\dd t}\norm{\chi v}_{L^2}^2
+\lambda\norm{\chi\nabla v}_{L^2}^2
\leq C\norm{v}_{L^2(\supp\chi)}^2.
\end{equation}
Integration over $[0,t]$ gives the $q=0$ case of \eqref{eq:energy}. Because $v$ belongs to the local Hilbert triple
\[
H^1(U)\subset L^2(U)\subset H^{-1}(U),
\]
its weak time derivative and the energy assumptions give an $L^2$-continuous representative. The calculation for weak solutions is justified by Steklov averages (or a time-mollification argument).

\subsection{Higher integer orders}
Let $1\leq q\leq k+1$ and $\abs{\alpha}\leq q$. Set $v_\alpha=D^\alpha v$. Differentiation formally gives
\begin{equation}\label{eq:diff}
(v_\alpha)_t+Pv_\alpha=F_\alpha,
\qquad F_\alpha=-[D^\alpha,P]v.
\end{equation}
The highest-order commutator is
\begin{equation}\label{eq:commutator}
[D^\alpha,a^{ij}\partial_i\partial_j]v
=\sum_{0<\beta\leq\alpha}\binom{\alpha}{\beta}
(D^\beta a^{ij})\,\partial_i\partial_jD^{\alpha-\beta}v.
\end{equation}
Every $\beta$ in this sum is nonzero; therefore
\[
2+\abs{\alpha-\beta}\leq q+1.
\]
Thus the commutator costs at most $q+1$ spatial derivatives of $v$, and it differentiates the coefficients at most $q$ times. Those derivatives are bounded since $q\leq k+1$.

Testing \eqref{eq:diff} against $\chi^2v_\alpha$, integrating by parts, and applying Young's inequality gives schematically
\begin{equation}\label{eq:highenergy}
\begin{aligned}
\frac{\dd}{\dd t}\norm{\chi v_\alpha}_{L^2}^2
+c\norm{\chi\nabla v_\alpha}_{L^2}^2
\leq C\norm{v}_{H^q(\supp\chi)}^2
+C\norm{v}_{H^q(\supp\nabla\chi)}^2.
\end{aligned}
\end{equation}
The terms of order $q+1$ arising from the commutator are absorbed into the coercive gradient term; lower-order errors remain on the right. Summing over $\abs\alpha\leq q$ and using nested spatial cutoffs yields \eqref{eq:energy} by induction on $q$ and Gronwall's inequality.

For continuity, the differentiated equation gives, locally,
\begin{equation}\label{eq:triple}
D^\alpha v\in L^2(0,T';H^1(U)),\qquad
\partial_tD^\alpha v\in L^2(0,T';H^{-1}(U)).
\end{equation}
The Hilbert-triple continuity theorem implies $D^\alpha v\in C([0,T'];L^2(U))$. Summing over $\abs\alpha\leq q$ yields $v\in C([0,T'];H^q(U))$.

To justify differentiating a weak solution, one first works with spatial difference quotients or regularized solutions, proves bounds independent of the regularization, and passes to the limit. The use of nested cutoffs is the same mechanism as in the interior elliptic argument of \cite{Taylor2023}: commutation with a cutoff lowers the order by one.

\section{Time differentiation and the elliptic bootstrap}\label{sec:time}

Because $P$ is independent of $t$,
\begin{equation}\label{eq:commtime}
[\partial_t,P]=0.
\end{equation}
Hence, formally, $w=u_t$ solves
\begin{equation}\label{eq:w}
w_t+Pw=0.
\end{equation}
When $u_0$ has two additional spatial derivatives, the natural initial trace for $w$ is
\begin{equation}\label{eq:w0}
w(0)=-Pu_0.
\end{equation}
The following standard time-difference-quotient lemma makes this precise.

\begin{lemma}[Initial trace of the time derivative]\label{lem:time}
Let $0\leq q\leq k+1$ be an integer. If $u_0\in H^{q+2}_{\loc}$, then $w=u_t$ is a local energy solution up to $t=0$ of \eqref{eq:w}, with initial trace $-Pu_0\in H^q_{\loc}$. Consequently,
\begin{equation}\label{eq:ut-cont}
u_t\in C([0,T);H^q_{\loc}).
\end{equation}
\end{lemma}

\begin{proof}[Explanation of the standard regularization argument]
For smooth solutions, test the equation against $\chi^2u_t$. Integration by parts gives
\begin{equation}\label{eq:test-ut}
\begin{aligned}
&\int\chi^2\abs{u_t}^2
+\frac12\frac{\dd}{\dd t}\int\chi^2a^{ij}\partial_iu\,\partial_ju\\
&\qquad\leq C\int_{\supp\chi}(\abs{u}^2+\abs{\nabla u}^2)
+\frac12\int\chi^2\abs{u_t}^2.
\end{aligned}
\end{equation}
No $\partial_ta^{ij}$ term occurs. Together with the preceding local energy bound this yields $u_t\in L^2_{\loc}([0,T)\times M^\circ)$.

To identify the initial value rigorously, extend $u$ by $u_0$ for negative times:
\[
\widetilde u(t)=\begin{cases}u_0,&t<0,\\u(t),&t\geq0.\end{cases}
\]
Continuity of the $L^2$ trace prevents a Dirac mass in $\partial_t\widetilde u$. Distributionally,
\begin{equation}\label{eq:extended}
(\partial_t+P)\widetilde u=\mathbf1_{\{t<0\}}Pu_0.
\end{equation}
For the forward difference quotient
\[
W_h(t)=\frac{\widetilde u(t+h)-\widetilde u(t)}{h},
\]
we obtain
\begin{equation}\label{eq:wh}
(\partial_t+P)W_h=-h^{-1}\mathbf1_{(-h,0)}Pu_0.
\end{equation}
The source on the right has uniformly bounded $L^1_tH^q_x$ norm on compact spatial regions:
\[
\int_{-h}^0\norm{h^{-1}Pu_0}_{H^q(U)}\,\dd t
=\norm{Pu_0}_{H^q(U)}.
\]
The inhomogeneous local parabolic estimate, combined with the established $L^2$ control of $u_t$, controls the regularized difference quotients. Passing to the limit identifies the jump of $W_h$ at $t=0$ with $-Pu_0$, gives the energy regularity of $w$, and yields \eqref{eq:ut-cont} by Lemma~\ref{lem:energy}.

\emph{Analytic input:} the preceding passage uses the standard inhomogeneous parabolic estimate with initial trace and time-difference-quotient approximation. It is a compact description of that standard argument, not a claim that the displayed $L^1_t$ bound alone proves the required $L^2_tH^{q+1}_x$ bound. A complete construction of the approximation and trace spaces is part of parabolic $L^2$ regularity theory; compare \cite{Krylov2008}.
\end{proof}

\subsection{Strong continuity at integer orders through $k+3$}
For integer $0\leq s\leq k+1$, Lemma~\ref{lem:energy} immediately gives
\[
u_0\in H^s_{\loc}\quad\Longrightarrow\quad
u\in C([0,T);H^s_{\loc}).
\]
Now take an integer $2\leq s\leq k+3$ and write $q=s-2$. Then $0\leq q\leq k+1$ and, if $u_0\in H^s_{\loc}$, Lemma~\ref{lem:time} gives
\begin{equation}\label{eq:ut-limit}
u_t(t)\longrightarrow-Pu_0
\quad\text{in }H^{s-2}_{\loc}\quad(t\downarrow0).
\end{equation}
Set $z(t)=u(t)-u_0$. From the equation,
\[
Pz(t)=-u_t(t)-Pu_0.
\]
The finite-regularity version of the local elliptic estimate discussed in \cite{Taylor2023} gives, for $U\Subset V\Subset M^\circ$,
\begin{equation}\label{eq:ellipticestimate}
\norm{z(t)}_{H^s(U)}
\leq C\br{\norm{Pz(t)}_{H^{s-2}(V)}+\norm{z(t)}_{L^2(V)}}.
\end{equation}
Both terms tend to zero: the first by \eqref{eq:ut-limit}, the second by initial $L^2$ continuity. Thus
\[
u(t)\longrightarrow u_0\quad\text{in }H^s(U).
\]
This reaches $s=k+3$ at integer orders. The final two spatial derivatives come from \emph{elliptic regularity}, rather than further differentiating the parabolic energy identity.

\section{Positive-time smoothing}\label{sec:smoothing}

Fix $0<\delta<T'<T$. Let $\eta(t)$ vanish near $\delta/2$ and equal one on $[\delta,T']$. Applying the positive-time version of the energy argument to $\eta u$ and iterating through the orders $0,1,\ldots,k+1$ gives
\begin{equation}\label{eq:pos1}
u\in C([\delta,T'];H^{k+1}_{\loc}).
\end{equation}
The temporal cutoff produces a source $\eta'(t)u$ supported away from $t=0$, and successive estimates use nested space-time cylinders.

Testing by localized $u_t$, now with a temporal cutoff, yields $u_t\in L^2_{\loc}(M^\circ\times(0,T))$. Time difference quotients show that $w=u_t$ is a local energy solution of the same homogeneous equation on positive-time cylinders. Applying the positive-time estimate to $w$ gives
\begin{equation}\label{eq:pos-ut}
u_t\in C((0,T);H^{k+1}_{\loc}).
\end{equation}
Since $Pu=-u_t$, the local elliptic estimate yields
\begin{equation}\label{eq:pos-elliptic}
\norm{u(t)}_{H^{k+3}(U)}
\leq C\br{\norm{u_t(t)}_{H^{k+1}(V)}+\norm{u(t)}_{L^2(V)}}.
\end{equation}
Therefore $u(t)\in H^{k+3}_{\loc}$ for every $t>0$. To see continuity in that space, fix $t_0>0$ and apply the elliptic estimate to $u(t)-u(t_0)$:
\begin{equation}\label{eq:pos-cont}
\begin{aligned}
\norm{u(t)-u(t_0)}_{H^{k+3}(U)}
\leq C\bigl(&\norm{u_t(t)-u_t(t_0)}_{H^{k+1}(V)}\\
&+\norm{u(t)-u(t_0)}_{L^2(V)}\bigr).
\end{aligned}
\end{equation}
Both terms tend to zero. Hence
\begin{equation}\label{eq:smoothingconclusion}
\boxed{u\in C((0,T);H^{k+3}_{\loc}(M^\circ)).}
\end{equation}
The mechanism is
\bluebox{$\underbrace{u_t\in H^{k+1}_{\loc}}_{\text{parabolic estimate}}
\quad\Longrightarrow\quad
\underbrace{u\in H^{k+3}_{\loc}}_{\text{elliptic estimate}}$.}

\section{Real, noninteger Sobolev exponents}\label{sec:fractional}

The preceding calculations show the result explicitly at integer orders. For real $0\leq s\leq k+3$, one uses fractional Sobolev estimates or a localized parabolic parametrix. In Euclidean coordinates define
\[
\Lambda^q=(I-\Delta)^{q/2}.
\]
For $0<q\leq k+1$, the commutator appearing in the equation for $\Lambda^qv$ is
\[
[\Lambda^q,a^{ij}]\partial_i\partial_jv.
\]
The relevant fractional commutator bound, for functions localized in nested coordinate neighborhoods, takes the schematic form
\begin{equation}\label{eq:fractionalcomm}
\norm{[\Lambda^q,a^{ij}]\partial_i\partial_jv}_{L^2}
\leq C\norm{v}_{H^{q+1}},
\end{equation}
with analogous bounds for lower-order coefficients. The crucial derivative count is again one derivative less than the principal term. The nonlocality of $\Lambda^q$ means that commutators with spatial cutoffs must also be treated; they are controlled on slightly larger regions and have lower-order or smoothing contributions.

The fractional localized energy estimate yields
\begin{equation}\label{eq:fractionalenergy}
L^2_tH^{q+1}_x\cap H^1_tH^{q-1}_x
\hookrightarrow C_tH^q_x,
\end{equation}
where all spaces are understood on nested interior space-time cylinders. This covers $0\leq q\leq k+1$; orders $q+2\leq k+3$ follow from the equation for $u_t$ and the corresponding fractional elliptic estimate. When $k=0$, the gap between $H^1$ and $H^2$ is handled by interpolation of the appropriate \emph{linear solution estimates}, not merely by interpolating an already unknown solution norm.

These are standard parts of the $L^2$ parabolic Sobolev theory, treated in greater detail in \cite{Krylov2008}. With this analytic input, the integer-order conclusions extend to all real $0\leq s\leq k+3$.

\begin{caution}[Why simple interpolation is not enough]
Suppose the datum lies only in $H^s$ for a noninteger $s$. Estimates proved for data in a strictly higher integer-order space cannot simply be interpolated for that particular datum: the higher norm might be infinite. One interpolates the solution maps or proves a fractional-order energy estimate and then uses density in the correct $H^s$ topology. This is precisely the additional step needed for the real-exponent formulation of Theorem~\ref{thm:main}.
\end{caution}

\section{Norm convergence and divergence}\label{sec:dichotomy}

\subsection{Finite limit when regularity holds beyond $\overline U$}
If $u_0\in H^r(V)$ for some $U\Subset V\Subset M^\circ$, initial-time continuity at order $r$ gives $u(t)\to u_0$ in $H^r(U)$. By the reverse triangle inequality,
\[
\abs{\norm{u(t)}_{H^r(U)}-\norm{u_0}_{H^r(U)}}
\leq\norm{u(t)-u_0}_{H^r(U)}\longrightarrow0.
\]
Therefore
\begin{equation}\label{eq:dich-finite}
\boxed{\lim_{t\downarrow0}\norm{u(t)}_{H^r(U)}=\norm{u_0}_{H^r(U)}.}
\end{equation}

\subsection{Infinite limit when the initial datum is not in $H^r(U)$}
Suppose $u_0\notin H^r(U)$. We know that $u(t)\to u_0$ in $\cD'(U)$ and that $u(t)\in H^r(U)$ for $t>0$. If $\norm{u(t)}_{H^r(U)}$ did not tend to infinity, there would exist $t_j\downarrow0$ with
\[
\sup_j\norm{u(t_j)}_{H^r(U)}<\infty.
\]
By weak compactness in the Hilbert space $H^r(U)$, a subsequence would converge weakly to some $v\in H^r(U)$. Weak $H^r$ convergence implies distributional convergence. Uniqueness of distributional limits then forces $v=u_0$, contradicting $u_0\notin H^r(U)$. Hence
\begin{equation}\label{eq:dich-inf}
\boxed{\lim_{t\downarrow0}\norm{u(t)}_{H^r(U)}=+\infty.}
\end{equation}
This argument is abstract: it applies to any family of $H^r(U)$ functions converging to $u_0$ in distributions, not only to the heat equation.

\section{Examples and counterexamples}\label{sec:examples}

\subsection{A boundary layer inside the chosen interior region}
\begin{example}[Regularity on $U$ alone does not give a finite limit]
Consider the classical heat equation on $\R$ with
\[
u_0(x)=\mathbf1_{(0,\infty)}(x).
\]
The solution is
\[
u(t,x)=\frac12\left(1+\operatorname{erf}\br{\frac{x}{2\sqrt t}}\right).
\]
Choose $U=(0,1)$. There $u_0\equiv1$, so $u_0|_U\in H^r(U)$ for every $r\geq0$. However,
\[
\partial_xu(t,x)=\frac1{\sqrt{4\pi t}}e^{-x^2/(4t)},
\]
and hence
\begin{equation}\label{eq:layer}
\begin{aligned}
\norm{\partial_xu(t)}_{L^2(0,1)}^2
&=\frac1{4\pi t}\int_0^1e^{-x^2/(2t)}\,\dd x\\
&\sim\frac1{4\sqrt{2\pi}}t^{-1/2}
\qquad(t\downarrow0).
\end{aligned}
\end{equation}
Consequently,
\[
\norm{u(t)}_{H^1(0,1)}\sim Ct^{-1/4}\longrightarrow+\infty.
\]
The initial jump lies at $x=0$, on the boundary of $U$; its diffusive layer extends a distance $O(\sqrt t)$ inside $U$. No neighborhood $V$ of $\overline U$ carries $H^1$ regularity of $u_0$. On every $U_\varepsilon=(\varepsilon,1-\varepsilon)$, by contrast, $u(t)\to1$ in $H^r(U_\varepsilon)$ for each fixed $r$.
\end{example}

\subsection{Sharpness of $H^{k+3}$ for $C^{k,1}$ coefficients}
\begin{example}[A persistent coefficient-induced singularity]
Let $\Omega=(-1,1)$ and
\[
a(x)=1+(x_+)^{k+1},\qquad x_+=\max\{x,0\}.
\]
Then $a>0$ and $a\in C^{k,1}(\Omega)$, but its $(k+1)$st derivative has a jump at $0$. Let
\[
P=-a(x)\frac{\dd^2}{\dd x^2}.
\]
Choose a positive first Dirichlet eigenfunction $\phi$:
\[
-a(x)\phi''(x)=\lambda\phi(x),\qquad
\phi(-1)=\phi(1)=0.
\]
In particular, $\phi(0)>0$. The eigenfunction equation reads
\begin{equation}\label{eq:eigen}
\phi''=-\lambda a^{-1}\phi.
\end{equation}
Since $a^{-1}\in C^{k,1}$, one-dimensional elliptic/ODE regularity gives $\phi\in H^{k+3}_{\loc}$. Differentiating \eqref{eq:eigen} $(k+1)$ times on either side of $0$ shows that the jump in the $(k+1)$st derivative of $a^{-1}$ produces a nonzero jump in $\phi^{(k+3)}$. Therefore $\phi\notin H^{k+4}_{\loc}$ near $0$.

Set
\[
u(t,x)=e^{-\lambda t}\phi(x).
\]
Then $u_t+Pu=0$, but for every $t>0$,
\[
u(t)\in H^{k+3}_{\loc}\setminus H^{k+4}_{\loc}.
\]
Hence $R=k+3$ cannot be raised to $k+4$ uniformly over $C^{k,1}$ coefficients. This does not assert that every such operator or every datum has precisely this regularity.
\end{example}

\begin{remark}[Fractional sharpness]
The simple jump example establishes sharpness against one \emph{full} extra derivative. To rule out a uniform gain of any prescribed fractional $\varepsilon>0$, one can choose rougher, oscillatory coefficients in $C^{k,1}$ whose top bounded weak derivative has inadequate fractional Sobolev regularity. This refinement is distinct from the explicit example above.
\end{remark}

\section{Corollaries for spatial derivatives}

\begin{corollary}[Strong convergence of weak derivatives]
If $u_0\in H^s_{\loc}$, $0\leq s\leq k+3$, then for every multi-index $\alpha$ with $\abs\alpha\leq s$,
\begin{equation}\label{eq:derivweak}
D^\alpha u(t)\longrightarrow D^\alpha u_0
\quad\text{in }H^{s-\abs\alpha}_{\loc}.
\end{equation}
In particular, $D^\alpha u(t)\to D^\alpha u_0$ strongly in $L^2(U)$ on each $U\Subset M^\circ$ whenever $\abs\alpha\leq s$.
\end{corollary}
\begin{proof}
Differentiation is bounded $H^s\to H^{s-\abs\alpha}$ on each coordinate neighborhood. Apply this map to initial-time strong convergence.
\end{proof}

\begin{corollary}[Uniform convergence of classical derivatives]
If $s>m+n/2$, where $m$ is a nonnegative integer, then for all $\abs\alpha\leq m$,
\begin{equation}\label{eq:derivstrong}
D^\alpha u(t)\longrightarrow D^\alpha u_0
\quad\text{uniformly on compact subsets as }t\downarrow0.
\end{equation}
At positive times one may guarantee continuous spatial derivatives of every integer order $m<k+3-n/2$ by Sobolev embedding alone.
\end{corollary}
\begin{proof}
On a slightly enlarged coordinate region the Sobolev embedding $H^s\hookrightarrow C^m$ is continuous if $s>m+n/2$. Apply it to $u(t)-u_0$, then use a finite coordinate cover.
\end{proof}

\section{What the qualitative theorem does not say}

\subsection{No universal blow-up rate}
If $u_0\in H^s_{\loc}\setminus H^r_{\loc}$, the theorem forces the $H^r(U)$ norm to diverge on each region where $u_0\notin H^r(U)$, but gives no universal power-law asymptotic. For the constant-coefficient heat equation, the Fourier multiplier gives, for $r\geq s$,
\begin{equation}\label{eq:rate}
\norm{e^{t\Delta}u_0}_{H^r(\R^n)}
\lesssim t^{-(r-s)/2}\norm{u_0}_{H^s(\R^n)},\qquad0<t\leq T,
\end{equation}
with constants allowed to depend on $T$. Variable-coefficient local analogues require a \emph{quantitative} parabolic smoothing theorem and typically a lower-order space-time norm on a larger region. The factor $t^{-(r-s)/2}$ is an upper bound, not a generic asymptotic equivalence. Additional initial regularity can give corresponding convergence rates by integrating $u_t=-Pu$.

\subsection{Regularity versus global well-posedness}
Interior regularity of a solution that already exists does not require specifying its boundary condition; existence, uniqueness, and quantitative control of that solution may require global hypotheses. Different Dirichlet, Neumann, Robin, or other admissible realizations can produce different interior solutions. They nevertheless obey the \emph{same local regularity estimates} provided they belong to the required local energy class.

\section{Proof strategy and comparison with elliptic regularity}

The argument is best understood as three consecutive mechanisms:
\begin{enumerate}
\item \textbf{Localized parabolic energy estimates:} cut off in space, differentiate (or use difference quotients), absorb commutators, and obtain $H^q$ estimates and strong time continuity for $q\leq k+1$.
\item \textbf{Time differentiation:} because $[\partial_t,P]=0$, the function $u_t$ solves the same parabolic equation. Its initial trace is $-Pu_0$ when the datum has two more derivatives.
\item \textbf{Localized elliptic bootstrapping:} the identity $Pu=-u_t$ and the local elliptic estimate yield the remaining two spatial derivatives, reaching $H^{k+3}_{\loc}$.
\end{enumerate}
The elementary ingredient behind the cutoff procedure is the commutator formula
\[
P(\chi v)=\chi Pv+[P,\chi]v,
\qquad\operatorname{ord}[P,\chi]\leq1.
\]
This is exactly the localization mechanism in the elliptic regularity results discussed in \cite{Taylor2023}. Taylor's development of elliptic a priori estimates, difference quotients, and localization supplies the elliptic half of the argument; for the fractional and initial-trace parabolic machinery, see \cite{Krylov2008}. The discussion here is a synthesis of those standard methods rather than a verbatim quotation or a separately named theorem from either book.

\section{Final consolidated conclusions}

For the time-independent, strongly elliptic, non-divergence-form operator \eqref{eq:operator} with $C^{k,1}_{\loc}$ coefficients and a local energy solution with initial trace $u_0$, one has:
\begin{center}
\small
\renewcommand{\arraystretch}{1.3}
\begin{tabularx}{\textwidth}{@{}>{\bfseries}p{0.26\linewidth}X@{}}
\toprule
Question & Conclusion\\
\midrule
Initial-time convergence & If $0\leq s\leq k+3$ and $u_0\in H^s_{\loc}$, then $u(t)\to u_0$ strongly in $H^s_{\loc}$.\\
Positive-time spatial smoothing & $u(t)\in H^{k+3}_{\loc}$ for every $t>0$, with continuity in that space on positive-time intervals; $k+3$ is the sharp general integer ceiling.\\
$H^r$-norm limit & For $0\leq r\leq k+3$, the norm on $U$ converges to the datum's norm if the datum is $H^r$ on a neighborhood of $\overline U$; it diverges if the datum is not in $H^r(U)$.\\
Boundary dependence & The local estimates do not involve the boundary condition, but global existence and bounds for the particular solution may.\\
\bottomrule
\end{tabularx}
\end{center}

\bluebox{$\displaystyle
u_t\in H^{k+1}_{\loc}
\quad\xRightarrow{\;Pu=-u_t\;}
\quad u\in H^{k+3}_{\loc},
\qquad
u_t(t)\to-Pu_0\ \Rightarrow\ u(t)\to u_0.
$}

\section*{References}
\addcontentsline{toc}{section}{References}
\begin{thebibliography}{9}
\bibitem{Taylor2023}
Michael E. Taylor, \emph{Partial Differential Equations I: Basic Theory}, third edition, Applied Mathematical Sciences, vol.~115, Springer, Cham, 2023. DOI: \href{https://doi.org/10.1007/978-3-031-33859-5}{10.1007/978-3-031-33859-5}. See especially the treatments of Sobolev spaces, linear elliptic equations, localized elliptic estimates, and difference quotients.
\bibitem{Krylov2008}
N. V. Krylov, \emph{Lectures on Elliptic and Parabolic Equations in Sobolev Spaces}, Graduate Studies in Mathematics, vol.~96, American Mathematical Society, Providence, RI, 2008. \href{https://doi.org/10.1090/gsm/096}{DOI: 10.1090/gsm/096}. Standard background for parabolic Sobolev estimates and initial-value regularity.
\end{thebibliography}

\end{document}
